% !TeX program = xelatex
\documentclass{blog}
\usepackage{fontspec}
\setmonofont{DejaVu Sans Mono}[Scale=MatchLowercase]

\newtheorem{tactic}{Tactic}[section]

\title{Lean}

\begin{document}
\maketitle
\tableofcontents
\section{Introduction}
\begin{definition}[lean]
    Lean is an expression-oriented functional language based on dependent type theory.
\end{definition}

\section{Dependent Type Theory}
\url{https://lean-lang.org/theorem_proving_in_lean4/}

Every expression has an associated type.

\begin{definition}[object]
\end{definition}

\begin{definition}[types]~
    \begin{verbatim}
Nat
Int
Float
String
Bool -- True or False
List
Type
\end{verbatim}
\end{definition}

\begin{definition}[declare object]~
    \begin{verbatim}
def obj : type := value -- define constant
#check obj -- check type
#eval expr -- evaluate
\end{verbatim}
\end{definition}

\begin{definition}[proposition]
\end{definition}

\begin{definition}[theorem]
\end{definition}

\begin{definition}[function]~
    \begin{verbatim}
def func (arg1 : type) ... (argn : type) : type :=
    body
\end{verbatim}
\end{definition}

\section{Tactics}
\begin{tactic}[reflexivity of equality]
    \verb|ref| proves goals of the form \verb|X = X|.
\end{tactic}

\begin{tactic}[rewrite]
    If \verb|h| is a proof of an equality \verb|X = Y|,
    then \verb|rw [h]| replace all \verb|X| with \verb|Y|.
    \begin{verbatim}
repeat rw [h] -- rewrite at once
rw [← h] -- turn Y into X
\end{verbatim}
\end{tactic}

\begin{definition}[$\mathbb{N}$]
    $\mathbb{N}$ is the natural numbers defined by \verb|0 : ℕ| and \verb|succ (0 : ℕ) : ℕ|.
\end{definition}

\begin{theorem}[successor]
    The successor of \verb|n| means the number after \verb|n|.
    \begin{verbatim}
two_eq_succ_one -- proof of 2 = succ 1
add_succ x d : x + succ d = succ (x + d)
\end{verbatim}
\end{theorem}

\begin{theorem}[adding 0]
    Define \verb|a + 0| to be \verb|a| for any number \verb|a|.
    \begin{verbatim}
add_zero   -- proof of ? + 0 = ?
add_zero x -- proof of x + 0 = x
\end{verbatim}
\end{theorem}

\appendix
\section{Tools}
\begin{definition}[comment]~
    \begin{verbatim}
/- comment blocks can be nested -/
-- line comment
\end{verbatim}
\end{definition}

% \printbibliography
\end{document}
